\documentclass[10pt,a4paper]{amsart}
\usepackage[margin=19mm]{geometry}
\usepackage{amsmath,amssymb,mathtools}
\usepackage[scr=rsfs,cal=euler]{mathalfa}
\usepackage{xfrac}
\usepackage[unicode,hidelinks,bookmarks=false]{hyperref}
\newcommand{\ie}{i.e.\ }
\newcommand{\cf}{cf.\ }
\newcommand{\resp}{respectively\ }
\newcommand{\Subsection}{\subsection}
\providecommand{\nobreakdash}{\nobreak\hbox{-}}
\setlength{\emergencystretch}{2em}
\multlinegap0pt
\usepackage{amssymb}
\usepackage{xcolor}
\newtheorem{theorem}{Theorem}
\title{Gradient systems and asymmetric relaxations in view of Riemannian geometry}
\author{Alessandro Bravetti, Miguel \'Angel Garc\'ia Ariza, Jos\'e Roberto Romero-Arias}
\date{Source: arXiv:2604.00192v1 (CC BY 4.0)}
\begin{document}
\maketitle

\noindent\emph{Source and licence.} Gradient systems and asymmetric relaxations in view of Riemannian geometry. Version arXiv:2604.00192v1 (CC BY 4.0), distributed by arXiv under the Creative Commons Attribution 4.0 International licence, http://creativecommons.org/licenses/by/4.0/, as recorded in the arXiv licence metadata for this exact version and shown on the abstract page https://arxiv.org/abs/2604.00192v1. Full licence notice: \texttt{licenses/arxiv-2604.00192-CC-BY-4.0.txt}.\par\smallskip

\noindent\textit{Editorial scope.} Counted unit: the article's theorem thm:asymmetric (source0.tex lines 320--331). The article's complete original proof of it is delivered with it (source0.tex lines 333--358, one \texttt{proof} environment of 26 lines). No mathematics is written, repaired or re-derived: the statement and the proof are the article's own lines, in the article's own notation, and no retained line is substituted or re-flowed.  Retained uncounted as the source-local context the proof needs: lines 190--192; lines 205--213.  Nothing was omitted from any retained span, so the package carries no omission record.  No retained line contains a citation, so the wrapper carries no bibliography and no bibliography entry.  No retained line contains a figure environment, an image-inclusion command or drawing code, so no image is delivered and the package declares no asset dependency.  Editorial formatting only, in addition: an AMS \texttt{amsart} preamble that restates the article's own declarations and macros used by the retained lines, namely the environments \texttt{theorem} on separate counters, together with the standard packages those lines need.  The retained environments print in this document as Theorem 1 in order of appearance, whereas the article prints the same items with the numbering of its own full document.  The complete native source of the article as distributed in the arXiv e-print for this exact version is bundled unchanged as \texttt{sources/197573/source0.tex}.
\begin{equation}\label{eq:gradientpregeodesic}
    \nabla^f_{\operatorname{grad}f}\operatorname{grad}f=\lambda \operatorname{grad }f,
\end{equation}

\begin{theorem}\label{thm:main}
    For any function $f$, there exists a symmetric connection $\tilde\nabla^f$ defined on $M$ minus the critical points of $f$ by
    \begin{equation}\label{eq:nablaf}
        \tilde\nabla^f_XY:=\nabla^g_XY- 
        g(X,Y)\frac{\nabla^g_{\operatorname{grad}f}\operatorname{grad}f-\lambda\operatorname{grad }f}{||\operatorname{grad }f||^2},
    \end{equation}
    for some real function $\lambda$ and satisfying Eq.~\eqref{eq:gradientpregeodesic}. 
    Hence, every gradient curve of $f$ is a pregeodesic of $\tilde\nabla^f$.
\end{theorem}

\begin{theorem}\label{thm:asymmetric}
     Let $f$ be a real function on $M$ attaining its unique
     minimum at $q$. Let $\nabla^f$ straighten $f$ with $\lambda\equiv\text{const.}$ 
     Let $\gamma_1$ and $\gamma_2$ be two integral curves of the gradient descent of $f$ 
     with initial conditions $f$-equidistant from $q$.
    If 
    \begin{equation}
        C^f(\dot\gamma_1,\dot\gamma_1,\dot\gamma_1)<C^f (\dot\gamma_2,\dot\gamma_2,\dot\gamma_2),
    \end{equation}
    whenever $||\dot\gamma_1||=||\dot\gamma_2||$,
    then $\gamma_1$ is faster than $\gamma_2$.
\end{theorem}

\begin{proof}
    We want to see that, with the above hypotheses, $\Delta f:=f(\gamma_2(t))-f(\gamma_1(t))\geq0$, for all $t\geq 0$. Given that $f$ is sufficiently well behaved, this would follow if all critical values of $\Delta f$ are maxima, since $\Delta f(0)=\lim_{t\to\infty}\Delta f=0$. 
    Denote by $\gamma$ an integral curve of the gradient descent of $f$. Then $\dot\gamma=-\operatorname{grad}f$ and
    \begin{equation}\label{eq:fdot}
        \dot f=\mathrm{d}f(\dot\gamma)=g(\operatorname{grad}f,\dot\gamma)=-||\dot\gamma||^2.
    \end{equation}
    So, when $||\dot\gamma_1||=||\dot\gamma_2||$, $\Delta f$ attains a critical value, whose nature is determined by the sign of $\Delta\ddot f.$
    From the last equation, we have that
    \begin{equation}
        \ddot f=-\dot\gamma[g(\dot\gamma,\dot\gamma)]=-\nabla^f_{\dot\gamma}g (\dot\gamma,\dot\gamma)-2g(\nabla^f_{\dot\gamma}\dot\gamma,\dot\gamma)=-C^f(\dot\gamma,\dot\gamma,\dot\gamma)-2\lambda g(\dot\gamma,\dot\gamma)=-C^f(\dot\gamma,\dot\gamma,\dot\gamma)-2\lambda \dot f,\label{eq:ddotf}
    \end{equation}
    using Eq.~\eqref{eq:gradientpregeodesic} in the third equality. 
    
    Theorem~\ref{thm:main} guarantees the existence of a connection $\nabla^f$ that straightens $f$ with $\lambda\equiv\text{const.}$ For such a connection, at a critical point $t_*$ we have
    \begin{equation}
    \Delta \ddot f(t_*)=-C^f(\dot\gamma_2, \dot\gamma_2,\dot\gamma_2)\vert_{t_*}+C^f(\dot\gamma_1,\dot\gamma_1,\dot\gamma_1)\vert_{t_*}.
    \end{equation}
    The last expression is negative if and only if
    \begin{equation}
        C^f(\dot\gamma_1,\dot\gamma_1,\dot\gamma_1)\vert_{t_*}<C^f(\dot\gamma_2,\dot\gamma_2,\dot\gamma_2)\vert_{t_*},\label{eq:main}
    \end{equation}
    and $\Delta f(t_*)$ is a maximum.
    
    Thus, if \eqref{eq:main} holds every time the descent speeds are equal, the relaxation along $\gamma_1$ will be faster than that along $\gamma_2$.

\end{proof}

\end{document}
